\subsection*{Висновки до розділу 4}

У четвертому розділі роботи описано методологію експериментального дослідження
автоматичного перекладу усного мовлення з української на англійську мову.
Основні результати розділу:

\begin{enumerate}

\item \textbf{Сформовано експериментальний корпус} із аудіозаписів
публічних звернень українською мовою загальною тривалістю близько 45 хвилин,
сегментованих на 12 фрагментів для оптимальної обробки моделями.
Корпус характеризується високою тематичною когерентністю та містить
типові для публіцистичного стилю мовні конструкції, імена власні та
абревіатури.

\item \textbf{Реалізовано автоматизований pipeline} обробки мовлення,
що поєднує модель автоматичного розпізнавання мовлення Whisper medium
та спеціалізовану модель машинного перекладу OPUS-MT.
Продуктивність системи (RTF = 0,26$\times$) підтверджує
її придатність для офлайн-обробки великих корпусів навіть на
обмежених обчислювальних ресурсах без використання GPU.

\item \textbf{Обґрунтовано вибір системи метрик} для оцінки якості
перекладу. Комбінація автоматичних метрик BLEU (лексичний збіг),
BERTScore (семантична подібність) та METEOR (морфологічний збіг із
урахуванням синонімів) дає повну кількісну оцінку.
Доповнення ручною анотацією помилок за типами (на основі фреймворку
MQM) та оцінкою міжфразової когерентності дає змогу виявити якісні
нюанси, недоступні для автоматичного аналізу.

\item \textbf{Розроблено класифікацію помилок перекладу} на основі
фреймворку MQM, адаптовану до специфіки дослідження. Виділено шість
основних типів помилок (пропуски, додавання, спотворення, помилки
в іменах власних, граматичні та стилістичні помилки) з трирівневою
шкалою серйозності.

\item \textbf{Описано дизайн експерименту,} що включає процедури
автоматичного та експертного оцінювання, статистичні методи аналізу
(описова статистика, кореляція Спірмена, критерій Манна-Уітні) та
заходи із забезпечення валідності дослідження.

\item \textbf{Передбачено аналіз каскадного ефекту} — впливу помилок
ASR на якість машинного перекладу через обчислення WER та його
кореляції з метриками якості перекладу.

\end{enumerate}

Представлена методологія створює основу для проведення експериментального
дослідження, результати якого наведено в наступному розділі.
